% ==================================================================
%  Correctness.tex  --  Section 3.2.4 "Correctness" (summary)
%  Loaded from main.tex.  The derivations, the proof of Theorem 1, sign
%  bootstrapping/cleaning and the error analysis are in
%  Appendix/CorrectnessDetails.tex (Appendix A).
% ==================================================================

\subsubsection{Correctness}
\label{sec:correct}

Every stage of the server rests on $\mathsf{Sgn}$, which is reliable only
outside its dead zone $(-2^{-\alpha}, 2^{-\alpha})$
(Section~\ref{sec:comparison}). A comparison is therefore correct when
the inputs that decide its output lie outside the dead zone. We fix here
the precision of each stage; Appendix~\ref{app:correct} gives the
derivations, the proof of Theorem~\ref{thm:sandwich}, and how the server
brings the outputs close to $\pm 1$.

\paragraph{\emph{Neighbor Check}}
Squared distances are continuous, so no precision keeps every input of
Eq.~\eqref{eq:cmpinput} out of the dead zone. The server uses
% Note: Normalize.py also clamps alpha_d from below at 15, which binds
%   only for n < 128.
$\alpha_d = \min(20,\, \lceil \log_2 n \rceil + 8)$ and moves the dead
zone entirely to the non-neighbor side by the threshold
\begin{equation}
  \label{eq:theta}
  \theta = \varepsilon_n^2 + \mu,
  \qquad \mu = 2^{-\alpha_d} B ,
\end{equation}
with $B$ from Eq.~\eqref{eq:bound}. No true neighbor is then missed; only
a non-neighbor with a squared distance between $\varepsilon_n^2$ and
$\varepsilon_n^2 + 2\mu$ may add a value between $0$ and $1$ to a count.
Writing $\mathrm{Nei}_r(i)$ for the number of points within distance $r$
of point~$i$, itself included, and
$\varepsilon_{\mathrm{hi}} = \sqrt{\varepsilon_n^2 + 2\mu}$, the count
$\hat{c}_i$ of Eq.~\eqref{eq:count} satisfies
\begin{equation}
  \label{eq:sandwich}
  \mathrm{Nei}_{\varepsilon_n}(i) \;\le\; \hat{c}_i \;\le\;
  \mathrm{Nei}_{\varepsilon_{\mathrm{hi}}}(i)
  \qquad \text{for every } i .
\end{equation}

\paragraph{\emph{Core Check}}
An integer count lies at least $\tfrac{1}{2}$ from
$\mathit{MinPts} - \tfrac{1}{2}$, so the input of
Eq.~\eqref{eq:coreinput} has magnitude at least $1/(2n)$, and
$\alpha_c = \lceil \log_2 (2n) \rceil$ keeps it outside the dead zone.

\begin{theorem}
\label{thm:sandwich}
Let $\mathrm{Core}_r$ denote the set of core points of plaintext DBSCAN
run at radius $r$ with the same $\mathit{MinPts}$, and
$\widehat{\mathrm{Core}}$ the set of points whose slot of
$\mathsf{Core\_mask}$ is not $0$, with $\alpha_c$ as above. Then
\[
  \mathrm{Core}_{\varepsilon_n} \;\subseteq\;
  \widehat{\mathrm{Core}} \;\subseteq\;
  \mathrm{Core}_{\varepsilon_{\mathrm{hi}}} .
\]
In particular $\widehat{\mathrm{Core}} = \mathrm{Core}_{\varepsilon_n}$
whenever no pair of points lies at a distance strictly between
$\varepsilon_n$ and $\varepsilon_{\mathrm{hi}}$.
\end{theorem}

\paragraph{\emph{Label Propagation}}
Two distinct labels differ by at least $1/n$, so after the division by
$1.1$ in Eq.~\eqref{eq:step}, $\alpha_\ell = \lceil \log_2 (1.1\,n) \rceil$
keeps the input of $\mathsf{Sgn}$ outside the dead zone. Two equal labels
give the input $0$, but $\mathsf{Max}$ then returns their common value
whatever the sign. The three precisions depend on $n$ alone, so the
public parameters fix them.
